\chapter{DNA chemistry and sequence geometry}
\label{ch:chemistry}

You receive two sequence files. One says \texttt{ACGTA}; the other says
\texttt{TACGT}. Are they complementary? A comparison of letters at the same
positions says no. A correctly oriented molecular drawing says yes. The
disagreement is not a biological mystery: the files and the drawing use
different coordinate conventions.

This chapter builds the missing translation layer. We start with the atoms
that give a DNA strand its direction, draw an antiparallel duplex, and then
implement the corresponding string operations. Our endpoint is a small
tested library for orientation, uncertainty symbols, and interval coordinates.
It predicts neither binding strength nor molecular structure. That boundary
is part of the implementation, not a disclaimer added after the calculation.

\section{Predict before you compute}

Write \texttt{ACGTA} from left to right, label its left end $5'$ and its
right end $3'$, and place a complementary strand beneath it. Predict:
\begin{enumerate}
\item the lower strand as it appears left to right in the drawing;
\item the lower strand recorded in the usual $5'$-to-$3'$ convention;
\item the original upper strand read backward without changing its bases.
\end{enumerate}
These are respectively \texttt{TGCAT}, \texttt{TACGT}, and \texttt{ATGCA}.
Only the second is the \Term{reverse complement}. Keeping these answers
separate prevents errors in alignment, tokenization, and annotation.
We need strings, zero-based indices, sets, and the idea of a covalent bond.
The chemistry explains why the string conventions exist.

\section{Where direction comes from}

\subsection{A nucleotide is more than a letter}

A \Term{nucleoside} contains a sugar attached to a nitrogenous base.
A \Term{nucleotide} also contains phosphate. In DNA the sugar is
2-deoxyribose. Prime marks distinguish its numbered atoms from atoms in the
base: $C1'$ belongs to the sugar, whereas an unprimed number refers to the
base's numbering system. The base attaches at $C1'$; the ring also contains
an oxygen, $O4'$. The $C5'$ carbon lies outside the five-membered sugar ring
\cite{iupac-conformation,iupac-symbols}.

\Plate{01-nucleotide}{Connectivity map of a DNA nucleotide. The sugar ring has four carbons and one oxygen; the fifth sugar carbon is outside it. The blue attachment is to a base, not another sugar. This is a topological projection: bond angles, stereochemistry, protonation, and most hydrogens are not specified.}{fig:nucleotide}

The deoxy distinction is local: DNA's $C2'$ has hydrogen where RNA's
corresponding position has a hydroxyl group. In unmodified DNA, $C2'$ is a
methylene carbon with two hydrogens. DNA does not lack oxygen everywhere:
the sugar ring, backbone, and some bases all contain oxygen.

The letter \texttt{A} records a base identity. It does not record sugar atom
coordinates, phosphate charge states, terminal phosphate occupancy, or
whether the molecule is free, paired, damaged, or chemically modified.
A sequence file is a useful projection, not a complete molecular specification.

\subsection{The phosphodiester backbone is an oriented path}

In an ordinary DNA backbone, a phosphate connects the $3'$ oxygen of one
sugar to the $5'$ oxygen associated with the next sugar. Follow the path:
\begin{equation}
C3'_i-O3'_i-P-O5'_{i+1}-C5'_{i+1}.
\label{eq:backbone}
\end{equation}
The two ester connections give the name \Term{phosphodiester}. This repeated
asymmetric connection makes the chain directional. Reading $5'$ to $3'$
is tied to chemical structure rather than to the left edge of a page.

\Plate{02-backbone}{Two nucleotide residues connected through a phosphodiester. The emphasized path is $C3'-O-P-O-C5'$; base attachments branch from the sugars. Nonbridging oxygen bonds use a conventional connectivity representation, not a unique localized electronic structure. Terminal groups cannot be inferred from sequence letters.}{fig:backbone}

A terminal $3'$ hydroxyl and a terminal $5'$ phosphate are common drawing
choices, not facts recoverable from every bare sequence string. A $5'$ end
can carry a hydroxyl or an engineered group. Likewise, a residue inside DNA
is not an isolated deoxynucleoside triphosphate: the latter has a different
phosphate inventory and is a substrate in polymerase-catalyzed synthesis.
Counting three phosphates per internal residue confuses precursor and product.

Our software deliberately omits this chemistry metadata. If an experiment
depends on terminal groups, store them explicitly beside the sequence.
Do not silently change what a base letter means.

\section{Pairing, stacking, and the limit of a string}

In a canonical Watson--Crick duplex, A pairs with T and G pairs with C.
The two strands run antiparallel \cite{nhgri}. A mechanism drawing should
distinguish covalent connectivity along each backbone, hydrogen-bond
interactions across a base pair, and stacking between neighboring layers.

\Plate{04-pairing-stacking}{Three interactions that a ladder drawing can blur. Solid side paths represent backbone connectivity; dashed cross-pair marks indicate pairing; shaded interlayer regions indicate stacking. Two A--T marks and three G--C marks recall canonical hydrogen-bond counts, but this plate does not resolve donor/acceptor atoms, helical twist, or free energies.}{fig:stacking}

It is tempting to assign two stability points to A--T and three to G--C,
add the points, and call the result a melting temperature. That calculation
has no temperature units and ignores neighboring bases, salt conditions,
strand concentration, and competing structures. It cannot deliver the
quantity it claims to predict.

Experimental work separating stacking and pairing contributions reinforces
why they should not be collapsed into a bond-count story
\cite{stacking,stacking-correction}. We use that work qualitatively:
no numerical energy decomposition or corrected figure is reproduced.
Chapter 6 will introduce a thermodynamic model with explicit assumptions.

Three questions must stay separate. Are the symbols complements? Can these
strands form this proposed duplex? What fraction is bound under specified
conditions? The first is a discrete relation; the second needs structural
and chemical assumptions; the third needs a physical model or measurement.
Passing the first test is not an answer to the third.

\Needspace{10\baselineskip}
\section{Build the reverse complement}

\subsection{The drawing and the stored sequence}

Let $s=s_0s_1\ldots s_{L-1}$ be stored $5'$ to $3'$. Define
\begin{equation}
c(A)=T,\quad c(T)=A,\quad c(C)=G,\quad c(G)=C.
\end{equation}
The aligned partner under position $i$ has identity $c(s_i)$, but its own
reading direction runs opposite to the upper strand. Its stored form is
\begin{equation}
\operatorname{rc}(s)_j=c(s_{L-1-j}),\qquad 0\leq j<L.
\label{eq:rc}
\end{equation}

\Plate{03-oriented-duplex}{One duplex, two reading directions. The lower row is \texttt{TGCAT} left to right, but \texttt{TACGT} along its own $5'$-to-$3'$ arrow. The separate stored-sequence line changes the representation, not the molecule: its arrow denotes neither synthesis nor a reaction.}{fig:duplex}

The complete function below is extracted from the tested reference file.
Its helper \texttt{complement} validates the alphabet and looks up each
symbol in a complement dictionary. Section~\ref{sec:ambiguity} derives that
dictionary from sets.
\Listing{reverse_complement}

\subsection{Properties are useful, but examples are indispensable}

Because $c(c(b))=b$, reversing and complementing twice restores every base:
\begin{equation}
\operatorname{rc}(\operatorname{rc}(s))=s.
\label{eq:involution}
\end{equation}
This is an \Term{involution}. It is a strong regression property but an
incomplete specification: the incorrect implementation
\texttt{return sequence} also passes. Our tests combine it with the
independently derived example \texttt{ACGTA} $\mapsto$ \texttt{TACGT}.

Concatenation exposes another relation:
\begin{equation}
\operatorname{rc}(xy)=\operatorname{rc}(y)\operatorname{rc}(x).
\label{eq:concat}
\end{equation}
Blocks change order as well as contents. Reverse-complementing each chunk
while preserving chunk order violates this equation. This will matter when
Evolutor groups bases into tokens.

Reversal preserves length; complementation swaps A/T counts and C/G counts.
GC fraction is unchanged. But several incorrect operations also preserve
GC fraction. Tests should target the information an operation must transform,
not just a convenient summary statistic.

\section{Ambiguity is a set, not a fifth base}
\label{sec:ambiguity}

An \Term{IUPAC ambiguity symbol} represents a set of possible identities.
R means A or G; Y means C or T; N permits all four. The official INSDC
table supplies the symbols \cite{insdc}. Complementation acts on the set:
\begin{equation}
C(S)=\{c(b):b\in S\}.
\end{equation}
Thus $C(\{A,G\})=\{T,C\}$, so R complements to Y. This follows from base
complementation rather than an unrelated second convention.

\begin{center}\input{\Generated/ambiguity.tex}\end{center}

The table is generated from the set definitions used by the tests. The
program builds a reverse dictionary from sets to symbols and derives every
ambiguity complement, avoiding two independent tables that could drift apart.

For aligned antiparallel symbols $a$ and $b$, define existential compatibility:
\begin{equation}
\mathrm{compatible}(a,b)
\quad\Longleftrightarrow\quad C(S_a)\cap S_b\ne\varnothing.
\label{eq:compat}
\end{equation}
R above Y passes because A--T and G--C are possible. R above R fails.
The first result does not mean every assignment pairs: A above C is an
allowed R/Y assignment but not a canonical complementary pair.

\Listing{compatible_aligned}

The lower string must be in \emph{aligned $3'$-to-$5'$ order}.
The docstring states this because the function name cannot carry the whole
orientation contract. Unequal lengths return false. Two empty strings return
true by the convention for an empty conjunction; this says nothing about
molecular association.

The expansion of \texttt{ARN} contains \texttt{AAA}, \texttt{AAC},
\texttt{AAG}, \texttt{AAT}, \texttt{AGA}, \texttt{AGC}, \texttt{AGG}, and
\texttt{AGT}: $1\cdot2\cdot4=8$ possibilities. No probability of $1/8$ has
been asserted. A set is not a distribution. Nor do independent ambiguity
symbols preserve correlations: if only \texttt{AA} and \texttt{GG} were
observed, writing \texttt{RR} also permits \texttt{AG} and \texttt{GA}.

\Listing{expansions}

The expansion budget is public behavior. Seven N symbols produce
$4^7=16{,}384$ strings, exceeding the default 4,096. Rejecting before
allocation makes the cost explicit. For large records, retain compact sets
or use bounded iteration instead of materializing every sequence.

We reject lowercase, RNA U, gaps, and whitespace. Other applications may
normalize them, but the policy must be declared. Lowercase can carry masking
information in some workflows; blindly uppercasing may discard an annotation.

\section{Reverse coordinates as well as letters}

An annotation $[a,b)$ selects indices $a$ through $b-1$ in a length-$L$
string. Reversal sends $i$ to $L-1-i$. The selected indices therefore run
from $L-b$ through $L-1-a$, giving
\begin{equation}
[a,b)\longmapsto[L-b,L-a).
\label{eq:interval}
\end{equation}
The endpoint $b$ was excluded. Treating both boundaries as included base
indices creates an off-by-one error.

\Plate{05-interval}{A length-nine annotation transformed with its strand. Ticks denote boundaries 0 through 9; bases occupy the spaces between ticks. Original \texttt{GTAG} in $[2,6)$ becomes \texttt{CTAC} in $[3,7)$ of the reverse complement. Both intervals contain four bases.}{fig:interval}

\Listing{reverse_interval}

For \texttt{ACGTAGCTA}, the reverse complement is \texttt{TAGCTACGT}.
The coordinate and string operations must satisfy
\begin{equation}
\operatorname{rc}(s)[L-b:L-a]=\operatorname{rc}(s[a:b]).
\label{eq:slice}
\end{equation}
Tests check every valid interval of a short ambiguity-containing sequence,
including empty intervals and end boundaries. Empty $[a,a)$ maps to empty
$[L-a,L-a)$; it must not acquire a base.

This changes local strand coordinates. It does not automatically convert
a genomic reference system, reorder transcript exons, or reinterpret
upstream/downstream labels. A record should declare its reference, strand,
coordinate origin, and endpoint convention. Correct formulas applied to
different reference systems can still produce an incorrect annotation.

\section{A palindrome is not a folding prediction}

A reverse-complement palindrome satisfies $s=\operatorname{rc}(s)$.
For canonical DNA of length $2m$, choosing the first $m$ bases forces the
remaining $m$ bases. Hence
\begin{equation}
P(L)=
\begin{cases}
4^{L/2},&L\text{ even},\\
0,&L\text{ odd}.
\end{cases}
\label{eq:palindrome}
\end{equation}
For odd length, the middle canonical base would have to complement itself;
none does. Length zero contributes one empty string by convention.

\Listing{palindrome_count}

Exhaustive enumeration through length six gives counts
$1,0,4,0,16,0,64$. The alphabet restriction matters: N is its own
\emph{set} complement, so the one-symbol string \texttt{N} is a symbolic
palindrome. This does not create a canonical base that pairs with itself.

A hairpin requires separated segments of one strand to form a stem with a
connecting loop. A string identity specifies neither loop geometry, steric
constraints, competing intermolecular duplexes, nor favorable free energy.
Self-complementarity is a sequence relation; formation of a specified
hairpin under specified conditions is a stronger claim.

\section{What the implementation establishes}

Our artifact validates an alphabet, computes reverse complements, maps
intervals, expands bounded uncertainty sets, and checks a combinatorial
result. It does not simulate atoms or infer hybridization. Golden results
in \path{drafts/ch05/results.json} are computed, not manually transcribed.

Run \texttt{python3 drafts/ch05/reference.py} for the examples and
\texttt{python3 -m pytest tests/test\_ch05\_reference.py} for the tests.
The PDF builder extracts whole tested functions and generates the table
from those same definitions.

The practical lesson is to separate the molecule, its oriented representation,
and the claims an algorithm supports. Failures then become local and
testable. Chapter 6 can add physical parameters without silently changing
what a sequence string means.

\Needspace{10\baselineskip}
\section{Exercises with worked solutions}

\subsection*{1. Three transformations}
For \texttt{AGTC}, write aligned complement, reversal, and reverse complement.
\textbf{Solution:} \texttt{TCAG}, \texttt{CTGA}, and \texttt{GACT}.
Only the final answer stores the partner $5'$ to $3'$.

\subsection*{2. Locate the oxygen}
A drawing labels five ring vertices as carbon and a branch as $C5'$.
Repair it. \textbf{Solution:} the ring has four carbons and $O4'$.
The branch gives five sugar carbons total. Atom identities matter more than
a visually pleasing pentagon.

\subsection*{3. Inspect a claimed backbone}
A diagram joins a base directly to the next phosphate and calls it the
backbone. \textbf{Solution:} the ordinary connection is
$C3'-O3'-P-O5'-C5'$. Bases branch from sugars at $C1'$; they are not
links in this backbone path.

\subsection*{4. Defeat an insufficient test}
Give a wrong function that passes the involution test.
\textbf{Solution:} identity and reversal-only both pass.
Require a known oriented duplex example as well as algebraic properties.

\subsection*{5. Reverse two blocks}
Let $x=\texttt{AC}$ and $y=\texttt{GTA}$. Check Eq.~\ref{eq:concat}.
\textbf{Solution:} $\operatorname{rc}(xy)=\texttt{TACGT}$;
$\operatorname{rc}(y)=\texttt{TAC}$, $\operatorname{rc}(x)=\texttt{GT}$.
Keeping the original block order would give the incorrect \texttt{GTTAC}.

\subsection*{6. Complement a set}
Why does B complement to V? \textbf{Solution:} $\{C,G,T\}$ becomes
$\{G,C,A\}=\{A,C,G\}$, represented by V. Set order is irrelevant.

\subsection*{7. Possible versus guaranteed}
Does aligned N/N always match? \textbf{Solution:} a complementary assignment
exists, so the predicate returns true. A/A is also allowed and does not
match. A probability requires a distribution over assignments.

\subsection*{8. Preserve correlations}
Encode the observations \texttt{AC} and \texttt{GT} with per-position
ambiguity. \textbf{Solution:} \texttt{RY} also allows \texttt{AT} and
\texttt{GC}, neither observed. Keep a joint representation if the correlation matters.

\subsection*{9. Reverse an annotation}
Map $[1,4)$ in length ten. \textbf{Solution:} $[6,9)$, still three bases.
Mapping again restores $[1,4)$. The tempting $[5,8)$ transforms boundaries
as if they were included indices.

\subsection*{10. Count and qualify}
How many length-eight canonical reverse-complement palindromes exist?
\textbf{Solution:} $4^4=256$. This counts strings, not stable hairpins.
Structure and conditions remain unspecified.

\subsection*{11. Bound the expansion}
How many strings does \texttt{NNRY} represent?
\textbf{Solution:} $4\cdot4\cdot2\cdot2=64$, within the default budget.
Seven N symbols exceed the budget; reject before full allocation.

\subsection*{12. Design a record contract}
Specify a modified strand with an interval annotation.
\textbf{Worked rubric:} include sequence, alphabet policy, storage direction,
reference identity, interval convention, strand, terminal modifications,
and provenance. State which fields reversal transforms. Explain why a
sequence-only function must not invent transformed chemical metadata.
